Um die Gleichheit $\mathrm{O}(n, \mathbb{C}) \cap \mathrm{U}(n) = \mathrm{O}(n)$ zu zeigen, betrachten wir die Definitionen der beteiligten Gruppen. Die orthogonale Gruppe $\mathrm{O}(n, \mathbb{C})$ besteht aus allen $n \times n$ Matrizen $A$ mit komplexen Einträgen, die die Bedingung $A^T A = I$ erfüllen, wobei $A^T$ die Transponierte von $A$ und $I$ die Einheitsmatrix ist. Die unitäre Gruppe $\mathrm{U}(n)$ besteht aus allen $n \times n$ Matrizen $B$ mit komplexen Einträgen, die die Bedingung $B^* B = I$ erfüllen, wobei $B^*$ die konjugiert-transponierte Matrix von $B$ ist.

Eine Matrix $C$ liegt im Schnitt $\mathrm{O}(n, \mathbb{C}) \cap \mathrm{U}(n)$, wenn sie sowohl orthogonal als auch unitär ist, das heißt, wenn sie die Bedingungen $C^T C = I$ und $C^* C = I$ erfüllt. Da $C^* = \overline{C^T}$, folgt aus $C^* C = I$, dass $C^T C = I$ und $C^* C = I$ nur dann gleichzeitig erfüllt sind, wenn $C$ eine reelle Matrix ist. Somit ist $C \in \mathrm{O}(n)$, der Gruppe der reellen orthogonalen Matrizen. Daher gilt $\mathrm{O}(n, \mathbb{C}) \cap \mathrm{U}(n) = \mathrm{O}(n)$.

Für die zweite Gleichheit $\mathrm{SO}(n, \mathbb{C}) \cap \mathrm{SU}(n) = \mathrm{SO}(n)$ betrachten wir die spezielle orthogonale Gruppe $\mathrm{SO}(n, \mathbb{C})$, die aus Matrizen $A \in \mathrm{O}(n, \mathbb{C})$ mit $\det(A) = 1$ besteht, und die spezielle unitäre Gruppe $\mathrm{SU}(n)$, die aus Matrizen $B \in \mathrm{U}(n)$ mit $\det(B) = 1$ besteht.

Eine Matrix $C$ liegt im Schnitt $\mathrm{SO}(n, \mathbb{C}) \cap \mathrm{SU}(n)$, wenn sie sowohl in $\mathrm{SO}(n, \mathbb{C})$ als auch in $\mathrm{SU}(n)$ liegt, das heißt, wenn sie die Bedingungen $C^T C = I$, $C^* C = I$ und $\det(C) = 1$ erfüllt. Wie zuvor folgt aus $C^T C = I$ und $C^* C = I$, dass $C$ reell ist, und somit ist $C \in \mathrm{SO}(n)$. Daher gilt $\mathrm{SO}(n, \mathbb{C}) \cap \mathrm{SU}(n) = \mathrm{SO}(n)$.